\documentclass[11pt]{article}
\usepackage[margin=1in]{geometry}
\usepackage{amsmath,booktabs}
\title{Notes on wave motion}
\author{A Weftly sample}
\date{}
\begin{document}
\maketitle
\begin{abstract}
A simple model connects the shape of a wave to the way it moves.
This short example brings prose, equations, and a table together.
\end{abstract}
\section{A pattern in motion}
A wave carries a disturbance through space. Its shape changes with
position and time, while the underlying relationship stays the same.
For a travelling wave, we can write
\begin{equation}
  y(x,t)=A\sin(kx-\omega t).
\end{equation}
Here, $A$ is the amplitude, $k$ is the wave number, and $\omega$ is
the angular frequency. The speed follows from their ratio:
\[
  v=\frac{\omega}{k}=f\lambda.
\]
\section{Reading the relationship}
A longer wavelength produces a higher speed when the frequency is
held constant. The table uses illustrative values to show this
relationship.
\begin{center}
\begin{tabular}{lrr}
\toprule
Wave & Frequency (Hz) & Wavelength (m) \\
\midrule
A & 2.0 & 0.5 \\
B & 2.0 & 1.0 \\
C & 2.0 & 1.5 \\
\bottomrule
\end{tabular}
\end{center}
\section{The next question}
A useful model makes the next question easier to ask. What changes
when the medium changes, or when two waves meet?
\end{document}
